%% msheet-example.tex -- example for the msheet class and package.
%%
%% This file replaces the 2023 example, which still used the long-gone
%% "question"/"problem"/"msbluebox" API; that file is kept in attic/.
\documentclass[factsheet, white, rounded, onehalfspacing, 11pt, a4paper]{msheet}

\usepackage{mstuff}
\usepackage{msheet}
\usepackage[language=ngerman]{lipsum}

\title{An example showcasing the features of the msheet package}
\author{Marc Schlienger}

\attribution{Aufgaben nach \href{https://math.mit.edu/~gs/}{Gilbert Strang}.}
\license{\href{https://creativecommons.org/licenses/by/4.0}{CC BY 4.0}}

%% Everything above can also be said through the key/value interface, which
%% additionally replaces the old public \icon and \titlebackcolor commands:
%\msheetsetup{type=experiment, titlecolor=DarkCyan, logo=Rolf-Benz-Schule.png}

\newcommand{\commandToText}[1]{\texttt{\string#1}}

%%%%%%%%%%%%%%%%%%%%%%%%%%%%%%%%%%%%%%%%%%%%%%%%%%%%%%%%%%%%%%%%%%%%%%%%%%%%%%
\begin{document}

\maketitle

\section{The msheet class (msheet.cls)}

The icon on the left hand side of the title box is set according to the chosen
document type. Available options are \textit{\color{FireBrick}worksheet},
\textit{\color{FireBrick}factsheet}, \textit{\color{FireBrick}groupwork},
\textit{\color{FireBrick}schedulework} and \textit{\color{FireBrick}experiment}.
This document is of type \textit{\color{FireBrick}factsheet}.

We can force a rounded title box using the option
\textit{\color{FireBrick}rounded}. The background color of the title box might
be chosen as well. Available options are \textit{\color{DarkCyan}cyan},
\textit{\color{LightSlateGray}grey},
\colorbox{LightSlateGray}{\textit{\color{white}white}} and
\textit{\color{Yellow}yellow}.

The line spacing can be set via the options
\textit{\color{FireBrick}singlespacing},
\textit{\color{FireBrick}onehalfspacing} (this is the setting in this document)
and \textit{\color{FireBrick}doublespacing}.

All of these can also be set with \commandToText{\msheetsetup}, which works
in the preamble and in the document body:
\textit{\color{FireBrick}\commandToText{\msheetsetup\{type=experiment,
titlecolor=DarkCyan\}}}. The old public \commandToText{\icon} and
\commandToText{\titlebackcolor} commands no longer exist.

A license and an attribution might be given via
\textit{\color{FireBrick}\commandToText{\license\{\}}} and
\textit{\color{FireBrick}\commandToText{\attribution\{\}}} respectively (see
footer). Both are optional; if they are omitted, the corresponding part of the
footer is simply left out and a single warning is issued.

\section{The msheet style (msheet.sty)}

\subsection{Headings}

\commandToText{\mchapter} and \commandToText{\mtestpart} typeset centred
coloured headings:

\mchapter{Ein Kapitel}
\mtestpart{Ein Teil}

\subsection{Exercise headings}

\msection[\color{FireBrick} A subsection like task environment.]
Compute the average velocity between $t = 5$ and $t = 8$.
\begin{mtaskcolumns}(3)
	\task $f(t) = 6t$
	\task $f(t) = 6t+2$
	\task $f(t) = \frac{1}{2}at^2$
	\task $f(t) = t-t^2$
	\task $f(t) = 6$
\end{mtaskcolumns}

\msection*[\color{FireBrick} And the same as starred version.]
Compute the average velocity between $t = 5$ and $t = 8$.
\begin{mexccolumns}(3)
	\exc $f(t) = 6t$ \exc $f(t) = 6t+2$
	\exc $f(t) = \frac{1}{2}at^2$
	\exc $f(t) = t-t^2$
	\exc $f(t) = 6$
\end{mexccolumns}

\msectionr[\color{DodgerBlue} A paragraph like task environment.]
\begin{mexccolumns}(3)
	\exc $f(t) = 6t$ \exc $f(t) = 6t+2$
	\exc $f(t) = \frac{1}{2}at^2$
	\exc $f(t) = t-t^2$
	\exc $f(t) = 6$
\end{mexccolumns}

\msectionr*[\color{DodgerBlue} And the same as starred version.]
\begin{mtaskcolumns}(3)
	\task $f(t) = 6t$
	\task $f(t) = 6t+2$
	\task $f(t) = \frac{1}{2}at^2$
	\task $f(t) = t-t^2$
	\task $f(t) = 6$
\end{mtaskcolumns}

\medskip
\noindent The old names \commandToText{\msexercise} and
\commandToText{\mpexercise} still work; they are aliases of
\commandToText{\msection} and \commandToText{\msectionr}.

\mibnp{}

\subsection{Lists}

\begin{mtasklist}
	\item When you jump and fall back your height is $y = 2t - t^2$ in the right
		units.
	\begin{mtasklist}
		\item Graph this parabola and its slope.
		\item Find the time in the air and the maximum height.
		\item \label{ct} Prove: Half of the time you are above $y = \frac{3}{4}$.
			\mibgreen{tip 1}\mibgreen{solution 1}

			\mibblue{Basketball players hang in the air partly because of \ref{ct}.}
		\end{mtasklist}
	\item If the ball on the unit circle reaches $t$ degrees at the time $t$, find
		its position and speed and upward velocity. \mibred{recommended}
\end{mtasklist}

\begin{msectionlist}
	\item When you jump and fall back your height is $y = 2t - t^2$ in the right
		units.
	\begin{msectionlist}
		\item Graph this parabola and its slope.
	\begin{msectionlist}
		\item Graph this parabola and its slope.
		\end{msectionlist}
		\item Find the time in the air and the maximum height.
		\item \label{c} Prove: Half of the time you are above $y = \frac{3}{4}$.
			\mibgreen{tip 1}\mibgreen{solution 1}

			\mibblue{Basketball players hang in the air partly because of \ref{c}.}
		\end{msectionlist}
	\item If the ball on the unit circle reaches $t$ degrees at the time $t$, find
		its position and speed and upward velocity. \mibred{recommended}
\end{msectionlist}

\noindent A plain enumeration with the msheet labels:
\begin{menumerate}
	\item first
	\item second
	\begin{menumerate}
		\item nested
	\end{menumerate}
\end{menumerate}

\noindent Section-numbered enumeration (\texttt{menumeratex}):
\begin{menumeratex}
	\item first
	\item second
	\begin{menumeratex}
		\item nested
	\end{menumeratex}
\end{menumeratex}

\noindent And the ticked list:
\begin{mdone}(2)
	* Hausaufgabe abgegeben
	* Heft nachgeführt
\end{mdone}

\mibnp{}

\subsection{Boxes}

\begin{mtbred}[Theorem 1 -- l'H\^opitals rule]
	Suppose $f(x)$ and $g(x)$ both approach zero as $x \to a$. Then
	$\frac{f(x)}{g(x)}$ approaches the same limit as
	$\frac{f^\prime(x)}{g^\prime(x)}$, if that second limit exists:
	$$\lim_{x \to a}\frac{f(x)}{g(x)} = \lim_{x \to
	a}\frac{f^\prime(x)}{g^\prime(x)}.$$
\end{mtbred}

\begin{mtbgreen}[Experiment -- The gravitational force]
	\lipsum[1]
\end{mtbgreen}

\begin{mtbblue}[Example]
	\lipsum[1]
\end{mtbblue}

\begin{mtbcolor}[Gold]{A golden titled box}
	\lipsum[1]
\end{mtbcolor}

\begin{mbinfo}
	\lipsum[1]
\end{mbinfo}

\begin{mbred}
	An untitled red box.
\end{mbred}

\begin{mbgreen}
	An untitled green box.
\end{mbgreen}

\begin{mbblue}
	An untitled blue box.
\end{mbblue}

\begin{mbcolor}[Gold]
	An untitled box in an arbitrary \texttt{svgnames} colour.
\end{mbcolor}

\subsection{Inline boxes}

The inline boxes are meant for hints and markers inside running text:
\mibred{Achtung}, \mibgreen{Tipp}, \mibblue{Beispiel} and
\mibcolor[Gold]{beliebige Farbe}. \commandToText{\mibnp} puts a
\enquote{Bitte wenden!} marker at the foot of the page and breaks it.

\subsection{Scores}

\commandToText{\mscore} puts the score of an exercise into the margin and adds
it to a running total; \commandToText{\mtotalscore} prints that total.
\commandToText{\mresetscore} starts a new count.

\msection[Scored exercise]
Compute $\int_0^1 x^2\,\mathrm{d}x$. \mscore{4}

\msection[Another scored exercise]
Compute $\int_0^1 x^3\,\mathrm{d}x$. \mscore{2.5}

\mtotalscore

\medskip
\noindent \commandToText{\mresetscore} starts a new count, for a sheet with
several independently scored parts.
\mresetscore

\msection[A third exercise, counted separately]
Compute $\int_0^1 x^4\,\mathrm{d}x$. \mscore{3}

\mtotalscore

\mwish

\end{document}
